\documentclass[10pt,twocolumn]{article}
\usepackage[utf8]{inputenc}
\usepackage[margin=0.75in]{geometry}
\usepackage{amsmath,amssymb}
\usepackage{booktabs}
\usepackage{graphicx}
\usepackage{hyperref}
\usepackage{xcolor}
\usepackage{listings}
\usepackage{microtype}

\hypersetup{
    colorlinks=true,
    linkcolor=blue,
    citecolor=blue,
    urlcolor=blue
}

\lstset{
    basicstyle=\ttfamily\scriptsize,
    keywordstyle=\color{blue}\bfseries,
    commentstyle=\color{gray},
    stringstyle=\color{red},
    breaklines=true,
    frame=single,
    tabsize=2
}

\title{\textbf{CS2201 Assignment Report: Implementation and Evaluation of FLACK and FURBYS Micro-op Cache Replacement Policies}}
\author{
    \textbf{Group GN4 (Coding Assignment Track)} \\
    Abhishek Kumar (Roll: 250101001) \quad
    Kaustav Bhowal (Roll: 250101050) \quad
    Jayesh Bansal (Roll: 250101046) \\
    \textit{Assigned TA: Nikhil (n.khatuwal@iitg.ac.in)} \\
    Department of Computer Science and Engineering, IIT Guwahati
}
\date{November 08, 2026}

\begin{document}

\maketitle

\begin{abstract}
Modern x86 processors employ a dedicated micro-operation ($\mu$-op) cache to power down the legacy instruction cache and decode pipeline. In data center applications with massive instruction footprints, micro-op caches suffer from severe capacity and conflict thrashing under standard LRU replacement. Conventional theoretical optima such as Belady's MIN and flow-based offline optimal (FOO) fail because they ignore variable entry sizes, variable micro-op costs, partial hits via intermediate exit points, and asynchronous decode delays. In this work, we implement the offline near-optimal policy \textbf{FLACK} and the profile-guided online policy \textbf{FURBYS} modeled on the AMD Zen~3 microarchitecture (512 entries, 8-way set-associative). Our evaluation across data center workloads demonstrates that FLACK achieves a 30.38\%--34.44\% miss reduction over LRU (outperforming Belady by 2.3\%--2.7\%), while FURBYS provides a 9.76\%--16.87\% miss reduction with less than 1.0\% hardware metadata overhead (46 bits per set).
\end{abstract}

\section{Introduction and Architectural Motivation}
Variable-length CISC instruction decoding (such as x86) consumes up to 12.5\% of per-core power. The micro-op cache serves decoded micro-ops directly to the backend in 1 cycle, allowing clock-gating of the L1 instruction cache and the multi-stage decoder pipeline. 

However, in data center workloads (e.g., Clang, Cassandra, PostgreSQL), over 99\% of misses under LRU are capacity (88.3\%) or conflict (10.8\%) misses. Conventional policies fail due to three unique hardware realities:
\begin{enumerate}
    \item \textbf{Disproportionate Cost and Size:} A Prediction Window (PW) occupies variable entries ($s$) and contains 1 to $8\times s$ micro-ops ($c$). Evicting a 1-uop entry versus an 8-uop entry incurs the same object eviction in Belady, but radically different micro-op miss penalties.
    \item \textbf{Partial Hits and Overlapping Windows:} PWs sharing a start PC but terminating at different branch outcomes can partially hit at intermediate exit points without decoder involvement.
    \item \textbf{Asynchronous Decode Latency:} Lookups and insertions are out-of-order due to non-blocking accumulation buffers and a 5-cycle decoder latency.
\end{enumerate}

\section{FLACK and FURBYS Algorithms}

\subsection{FLACK (Offline Near-Optimal Solver)}
FLACK extends flow-based offline caching by:
\begin{equation}
\text{Eviction Score}_i = \frac{\text{Distance}_i}{\max\left(0.1, \frac{\text{cost}_i}{\text{size}_i}\right)}
\end{equation}
Lines with high micro-op density ($\text{cost}/\text{size}$) are retained longer. Furthermore, FLACK prioritizes keeping larger windows when future partial hits are predicted and buffers evictions lazily until insertions complete.

\subsection{FURBYS (Online Practical Policy)}
FURBYS translates FLACK insights into hardware using a 7-step profiling pipeline:
\begin{enumerate}
    \item \textbf{Trace Extraction:} Capture instruction streams via Intel PT.
    \item \textbf{Zero-Size Cache Run:} Record PW lookup sequences.
    \item \textbf{FLACK Offline Scoring:} Determine micro-op hit rates.
    \item \textbf{Jenks Natural Breaks:} 1D dynamic programming partitions per-set hit rates into 8 clusters (3-bit hints, $0\dots 7$) injected into branch instructions.
    \item \textbf{Hardware Execution:} Decoder forwards 3-bit hints to the accumulator and micro-op cache.
\end{enumerate}

\subsubsection{Hardware Metadata Overhead}
For an 8-way set:
\begin{equation}
\text{Overhead} = (3\text{ bits weight} + 2\text{ bits SRRIP}) \times 8 + 2 \times 3\text{ bits} = \mathbf{46\text{ bits}}
\end{equation}
Comparing against 4,608 bits of data per set, the overhead is only $\mathbf{0.998\%} < \mathbf{1.0\%}$.

\subsubsection{Replacement and Local Pitfall Detection Logic}
\begin{itemize}
    \item \textbf{Min Module:} Scans 8 ways to identify $W_{\min}$ and \texttt{min\_way} (tie-broken by max SRRIP RRPV).
    \item \textbf{Selective Bypassing:} If $W_{\text{new}} < W_{\min} - K$ (with $K=1$), insertion is bypassed, preventing cache pollution and saving write energy.
    \item \textbf{Local Miss-Pitfall Detector:} When \texttt{pitfall\_slot\_0 == min\_way} and \texttt{pitfall\_slot\_1 == min\_way}, the policy detects that a globally hot PW has become locally cold. It triggers fallback to 2-bit SRRIP (evicts $\text{RRPV} = 3$) and resets the pitfall registers.
\end{itemize}

\section{Evaluation Results}

Table~\ref{tab:results} summarizes the performance across evaluated policies on the AMD Zen 3 geometry.

\begin{table}[h]
\centering
\caption{Micro-op Cache Policy Comparison (Zen 3 Geometry)}
\label{tab:results}
\resizebox{\columnwidth}{!}{%
\begin{tabular}{lcccc}
\toprule
\textbf{Policy} & \textbf{Type} & \textbf{Uop Misses} & \textbf{Miss Rate} & \textbf{Red. vs LRU} \\
\midrule
LRU (Baseline) & Online & 122,949 & 100.00\% & 0.00\% \\
SRRIP (2-bit) & Online & 118,336 & 96.25\% & +3.75\% \\
Belady-MIN & Offline & 88,488 & 71.97\% & +28.03\% \\
\textbf{FLACK} & \textbf{Offline} & \textbf{85,594} & \textbf{69.62\%} & \textbf{+30.38\%} \\
\textbf{FURBYS (Same)} & \textbf{Online} & \textbf{109,898} & \textbf{89.38\%} & \textbf{+10.62\%} \\
FURBYS (Cross) & Online & 121,111 & 98.50\% & +2.14\% \\
\bottomrule
\end{tabular}%
}
\end{table}

\section{Sensitivity Analyses}
\begin{itemize}
    \item \textbf{Representative Bits:} 3 bits (8 groups) is the optimal balance point (+8.61\% miss reduction). Higher bits suffer diminishing returns and increase metadata storage.
    \item \textbf{Pitfall Detector Depth:} Depth 2 provides the highest stability (+8.61\% miss reduction, 1,183 fallbacks), preventing thrashing during phase shifts without over-triggering.
    \item \textbf{Selective Bypassing:} Bypasses ~28\% of cold insertions ($K=1$), shielding warm lines and saving dynamic write power.
\end{itemize}

\section{Conclusion}
We implemented and evaluated the complete FLACK and FURBYS micro-op cache replacement frameworks in C++ and gem5. FLACK sets a new theoretical reference bound (+30.38\% over LRU), while FURBYS delivers practical online miss reduction (+10.62\%) with $<1\%$ storage overhead.

\end{document}

